\begin{RequAnswer}{7.1}
We need a ``starting point.'' Induction proofs involve proving
a statement of the form $P(k)\rightarrow P(k+1)$. That is, we prove that if $P$ is
true for some value, then it is true for the next value.
But that does not imply that it is ever true--only that if it is true for some value,
it is true for the next value.
So we need to prove that $P$ is true for some value to get things started.
\end{RequAnswer}
\begin{RequAnswer}{7.2}
This is not circular reasoning because we are not proving that $P(k+1)$ is true.
We are proving that $P(k)\rightarrow P(k+1)$ is true.
In other words we are assuming $P(k)$ is true, and the using that fact to prove that $P(k+1)$,
which is a {\em different} statement, is true.
But again, we did not prove that $P(k+1)$
is {\em always} true. We only proved that {\em IF} $P(k)$ is true, then $P(k+1)$ is true.
\end{RequAnswer}
\begin{RequAnswer}{7.3}
(a) It is saying that if $P(a)$ is true and it is also true that whenever $P(k)$ is true that $P(k+1)$ is true,
then $P(n)$ is true for all $n\geq a$.
(b) We know that $P(a)$ is true.
We also know that $P(k)\rightarrow P(k+1)$ is true for all $k\geq a$.
In particular, we know that $P(a)\rightarrow P(a+1)$ is true.
Using modus ponens, we can conclude that $P(a+1)$ is true.
But then we can use modus ponens again to conclude that $P(a+2)$ is true
(since we know $P(a+1)$ and $P(a+1)\rightarrow P(a+2)$ are both true).
We can keep doing this over and over again so that eventually we can
show that $P(a+k)$ is true for any $k\geq 0$. Thus, $P(n)$ is true for all $n\geq a$.
\end{RequAnswer}
\begin{RequAnswer}{7.4}
No. Notice that we know that $P(0)$ is true, but we only know that $P(k)\rightarrow P(k+1)$ for $k>0$.
So we know $P(1)\rightarrow P(2)$, but we do not know whether or not $P(0)\rightarrow P(1)$.
In other words, we have a base case and an inductive case, but the inductive case does not go all the
way down to the base case, so we cannot connect them.
(Note: this is not a failure of induction. It is a failure in trying to use induction. A proper induction proof
would show that $P(k)\rightarrow P(k+1)$ for $k\geq 0$ so that the inductive step applies to the base case.)
\end{RequAnswer}
\begin{RequAnswer}{7.5}
There are clearly $2=2^1$ binary strings of length 1 (`0' and `1').
Assume that there are $2^k$ binary strings of length $k$.
Every string of length $k+1$ ends with either a `0' or a `1', and the first $k$ characters can
be any of the possible binary strings of length $k$. In other words, the number of binary strings of length
$k+1$ is $2\cdot 2^k=2^{k+1}$ since we can append to each of the $2^k$ strings of length $k$ either a `0' (producing
$2^k$ strings of length $k+1$) or a `1' (producing a different $2^k$ strings of length $k$).
Since we proved the base and inductive cases, we have shown that the number of binary strings of length $k$ is $2^k$.
\end{RequAnswer}
\begin{RequAnswer}{7.6}
These are both correct techniques.
\end{RequAnswer}
\begin{RequAnswer}{7.7}
In weak induction you assume $P$ is true for a given value and show $P$ is true for the next value.
For instance, you might assume $P(k)$ is true and prove that $P(k+1)$ is true.
In strong induction you assume that $P$ is true for every value from the base case up to a certain value,
and then you prove it for the next value.
For instance, you assume $P(1)\wedge P(2)\wedge\cdots \wedge P(k-1)$ is true
and prove that $P(k)$ is true.
(By the way, I purposely used $P(k)$ and $P(k+1)$ for one and $P(k-1)$ and $P(k)$ for the other
to re-emphasize that you can do it either way.)
\end{RequAnswer}
\begin{RequAnswer}{7.8}
Induction is like a bunch of dominoes lined up. If you push the first one over, it will push the next one over, which will push
the next one over, etc., until they have all fallen down. The first domino is the base case. The fact that the dominoes
are placed close enough to each other is like the inductive case (because they are close enough, if one falls, the next one will).
\end{RequAnswer}
\begin{RequAnswer}{7.9}
It links to itself, kind of like how a recursive algorithm calls itself.
\end{RequAnswer}
\begin{RequAnswer}{7.10}
(1) A recursive algorithm mush have one ore more {\em inductive cases} (or recursive case) where the algorithm calls itself.
This is required since otherwise the algorithm is not recursive.
(2) A recursive algorithm must have one or more {\em base cases} that can be solved directly (or at least non-recursively).
This is required since otherwise the algorithm would never finish.
(3) The recursive calls must be making progress toward the base cases.
This is also required since otherwise the algorithm would never finish.
\end{RequAnswer}
\begin{RequAnswer}{7.11}
They are both based on the same ideas, particularly that of bases cases and inductive cases.
\end{RequAnswer}
\begin{RequAnswer}{7.12}
Here is one solution.
This should be called with $n$ being the size of the array.
Notice that it searches from the end of the array to the beginning of the array
because it is the most straightforward way to implement the recursive idea.
\begin{code}
int search(int[] a, int n, int value) {
    if(n<=0) {
        return -1;
    } else if(a[n-1]==value) {
       return n-1;
    } else {
        return search(a,n-1,value);
    }
}
\end{code}
\end{RequAnswer}
\begin{RequAnswer}{7.13}
Neither is better. Iterative algorithms are often faster than their recursive equivalents,
but recursive algorithms are often easier to come up with and implement.
So they both have their merits.
\end{RequAnswer}
\begin{RequAnswer}{7.14}
A recurrence relation is a recursively defined formula for the values of a sequence.
In other words, it is a formula to compute $a_n$ based on one or more values of $a_i$
where $i<n$.
\end{RequAnswer}
\begin{RequAnswer}{7.15}
It means to come up with a closed formula. That is, formula to compute the $n$th term of the
sequence directly (not based on previous values of the sequence).
\end{RequAnswer}
\begin{RequAnswer}{7.16}
(a) The substitution method essentially involves guessing the correct formula and using induction to prove it.
(b) The iteration method keeps applying the recursive definition to the right side of the formula until it can
be simplified down to (generally speaking) a summation and the base case(s) that is then simplified.
(c) To use the Master Theorem, one verifying that the recurrence relation is in the correct form,
identifies the constants from the theorem ($a$, $b$, and $d$),
determines which case the formula falls into based on the values of the constants,
and writes down the answer based on which case it is.
This technique only gives an asymptotic bound on the solution, not an exact solution.
\end{RequAnswer}
\begin{RequAnswer}{7.17}
(a) They both give an exact formula whereas the Master Theorem only gives an asymptotic bound.
(b) The Master Method is {\em much easier} to use than the other two techniques.
(c) You need to be able to determine the answer before you prove it. If the formula is complicated,
you may not be able to determine what it is.
(d) The main downside is that it is messy. It also only works well on simple recurrence relations. For instance,
if a recurrence relation has several recursive terms (e.g. $T(n)=T(n-1)-2T(n-3)$), it would probably be
quite complicated to try to solve it using iteration.
(e) The Master Theorem also only works on one specific type of recurrence relation
and it does not give an exact solution.
(f) If I don't care about an exact solution, the Master Theorem is by far the easiest, so I would use that one.
If I want an exact solution, I would prefer substitution if I can see an obvious pattern and find the formula.
If I can't find a formula, I would prefer iteration because I should be able to work it out using that technique.
\end{RequAnswer}
\begin{RequAnswer}{7.18}
Because the three topics have a lot in common. Recurrence relations are just a form of recursion, and
they can be solved using induction.
\end{RequAnswer}
\begin{RequAnswer}{7.19}
Because recursive algorithms are often analyzed by developing and solving recurrence relations.
\end{RequAnswer}
\begin{RequAnswer}{7.20}
Develop a recurrence relation that describes the running time of the algorithm, including one or more base cases.
Use one of the techniques from the previous section to solve the recurrence relation.
\end{RequAnswer}
\begin{RequAnswer}{7.21}
Your answer may be different depending on your algorithm.
However, if your algorithm does not have a complexity of $\Theta(n)$,
you either have an incorrect algorithm (if its complexity is better than this),
a really bad algorithm (if its complexity is worse than this), or you analyzed it incorrectly.

Let $T(n)$ be the time it takes to run \verb+search(int[] a, int n, int value)+.
If $n\leq 0$, the algorithm just returns -1 which takes constant time. So $T(0)=1$.
If $n>=0$, the algorithm checks one value of the array, which takes constant time,
and then either returns or (in the worst case) makes a recursive call on an array of size $n-1$.
(technically the array still has size $n$, but we are telling the algorithm that it only has size $n-1$.)
So in this case, $T(n)=T(n-1)+1$.
This one is easy to solve by a variety of techniques. Let's do it using iteration.
First, notice that if $T(n)=T(n-1)+1$, then if we plug in $n-1$ for $n$, we would get $T(n-1)=T(n-2)+1$.
Therefore, $T(n)=T(n-1)+1=(T(n-2)+1) + 1)= T(n-2)+2$.
But $T(n-2)=T(n-3)+1$, so $T(n)=T(n-2)+2=(T(n-3)+1)+2=T(n-3)+3$.
It's not difficult to see that we can generalize this to $T(n)=T(n-k)+k$.
When $k=n$, we get $T(n)=T(n-n)+n= T(0)+n = n+1$.
So the worst-case performance of the algorithm on an array of size $n$ is $n+1$ steps.
\end{RequAnswer}
